\section{Baseline Benchmark Results}\label{sec:results}

This section establishes the empirical reference frame against which the
proposed model is judged. Ten architectures --- the nine standard baselines and
the proposed Custom EfficientNet~V2 --- are trained from scratch under the
single matched recipe of Section~\ref{sec:methods}, with no per-model
hyper-parameter tuning, and evaluated once on the held-out test set. The
benchmark is reported along three axes: classification accuracy on both heads
(Section~\ref{subsec:quant}), per-class subtype behaviour
(Section~\ref{subsec:perclass}), and computational efficiency
(Section~\ref{subsec:efficiency}). The ablation that explains \emph{why} the
proposed hub is designed as it is follows in Section~\ref{sec:ablation}; the
present section is concerned only with where the proposed model sits relative
to the field.

\subsection{Quantitative Comparison}\label{subsec:quant}

Table~\ref{tab:benchmark} reports the held-out test performance of all ten
architectures on both the binary (benign versus malignant) and the seven-class
subtype head, alongside parameter count, compute cost in GFLOPs,
and on-disk model size. Every value derives from a single training run per
model under the matched recipe; no figure in the table benefits from per-model
tuning, so the comparison isolates architecture rather than optimisation
effort.

\begin{table*}[!t]
\caption{Test-set classification performance under the matched fair recipe.
All ten architectures are trained from scratch (no ImageNet pre-training) with
identical optimiser, schedule, batch size, and seed. Bold marks the
proposed-model row and the best value in each column.}
\label{tab:benchmark}
\centering
\resizebox{\textwidth}{!}{%
\begin{tabular}{l d{4} d{4} d{4} d{4} d{2} d{3} d{2}}
\hline
\multicolumn{1}{l}{Model} &
\multicolumn{1}{c}{Binary Acc} &
\multicolumn{1}{c}{Binary F1} &
\multicolumn{1}{c}{Subtype Acc} &
\multicolumn{1}{c}{Subtype F1} &
\multicolumn{1}{c}{Params (M)} &
\multicolumn{1}{c}{GFLOPs} &
\multicolumn{1}{c}{Size (MB)} \\
\hline
ResNet50            & 0.9912 & 0.9912 & 0.9909 & 0.9909 & 25.61 & 4.134 & 98.01 \\
DenseNet121         & 0.9889 & 0.9889 & 0.9933 & 0.9933 &  7.92 & 2.834 & 31.15 \\
ConvNeXt-Tiny       & 0.9532 & 0.9533 & 0.9046 & 0.9048 & 28.59 & 4.456 & 109.22 \\
Swin-T              & 0.9819 & 0.9819 & 0.9769 & 0.9769 & 28.29 & 4.372 & 108.07 \\
EfficientNet-B0     & 0.9883 & 0.9883 & 0.9903 & 0.9903 &  5.28 & 0.386 & 20.60 \\
EfficientNet V2-B2  & \multicolumn{1}{c}{\textbf{0.9936}} & \multicolumn{1}{c}{\textbf{0.9936}} & 0.9933 & 0.9933 & 10.00 & 1.100 & 39.17 \\
EfficientNet V2-B3  & 0.9889 & 0.9889 & 0.9842 & 0.9842 & 14.21 & 1.522 & 55.57 \\
EfficientNet V2-S   & 0.9866 & 0.9865 & 0.9836 & 0.9836 & 21.22 & 2.711 & 82.86 \\
Inception V3        & 0.9930 & 0.9930 & 0.9921 & 0.9921 & 23.85 & 2.838 & 91.47 \\
\hline
\multicolumn{1}{l}{\textbf{Custom EfficientNet V2 (Hub v1, full)}} &
0.9906 & 0.9906 & 0.9921 & 0.9921 & 4.80 & 0.495 & 18.99 \\
\multicolumn{1}{l}{\textbf{Custom EfficientNet V2 (Hub v2, proposed)}} &
\multicolumn{1}{c}{\textbf{0.9906}} &
\multicolumn{1}{c}{\textbf{0.9907}} &
\multicolumn{1}{c}{\textbf{0.9951}} &
\multicolumn{1}{c}{\textbf{0.9951}} &
\multicolumn{1}{c}{\textbf{4.79}} &
\multicolumn{1}{c}{\textbf{0.493}} &
\multicolumn{1}{c}{\textbf{18.94}} \\
\hline
\end{tabular}%
}
\end{table*}

Two readings follow immediately from Table~\ref{tab:benchmark}. \textbf{First},
the proposed model attains the highest subtype accuracy in this benchmark. At
$99.51\,\%$ it exceeds the strongest baseline, EfficientNet~V2-B2 at
$99.33\,\%$, by $0.18$ percentage points, and the heaviest comparator,
Inception~V3 at $99.21\,\%$, by $0.30$ percentage points. \textbf{Second}, this
subtype lead is obtained while the proposed model is simultaneously the
smallest and cheapest architecture in the comparison. At $4.79$~M parameters it
is $2.1\times$ smaller than EfficientNet~V2-B2 ($10.0$~M), $5.0\times$ smaller
than Inception~V3 ($23.9$~M), and $5.3\times$ smaller than ResNet50
($25.6$~M), with proportional reductions in GFLOPs. The proposed model is
therefore not trading capacity for accuracy: among the nine baselines trained
from scratch under the matched recipe, it strictly dominates every baseline in
the subtype-accuracy versus parameter-count plane --- simultaneously the most
accurate and the smallest.

Figure~\ref{fig:pareto} makes this dominance explicit. Plotting subtype
accuracy against parameter count on a logarithmic axis, the proposed model
occupies the upper-left corner --- highest accuracy at the lowest parameter
budget --- and lies on the Pareto frontier, while every heavier comparator sits
strictly below and to the right of it. Figure~\ref{fig:radar} extends the
comparison to six axes, contrasting the proposed model against the two
strongest baselines on binary accuracy, subtype accuracy, and four
efficiency-derived axes (parameter, FLOPs, size, and memory efficiency, each
inverted so that a larger radius is always better). The proposed model encloses
the largest area: it matches the baselines on the two accuracy axes while
extending far beyond them on every efficiency axis.

\begin{figure}[!t]
\centering
\includegraphics[width=\columnwidth]{fig03_pareto_params}
\caption{Subtype accuracy versus parameter count (logarithmic horizontal axis)
for all ten architectures. The proposed Custom EfficientNet~V2 (Hub~v2) lies on
the Pareto frontier in the upper-left region --- the highest subtype accuracy
at the smallest parameter budget. ConvNeXt-Tiny ($90.5\,\%$) falls below the
plotted accuracy range and is marked separately.}
\label{fig:pareto}
\end{figure}

\begin{figure}[!t]
\centering
\includegraphics[width=\columnwidth]{fig03d_proposed_vs_baselines_radar}
\caption{Six-axis comparison of the proposed model against the two strongest
baselines, EfficientNet~V2-B2 and Inception~V3. The four efficiency axes are
inverted ($1/\text{metric}$) so that a point further from the centre is always
better. The proposed model matches the baselines on accuracy while dominating
on parameter, FLOPs, size, and memory efficiency.}
\label{fig:radar}
\end{figure}

The binary head, by contrast, is effectively tied at the top of the table.
EfficientNet~V2-B2 leads at $99.36\,\%$ and the proposed model follows at
$99.06\,\%$. The binary head is evaluated on a $1\,711$-image test set --- the
DS1 test split together with the full DS2 test set, each image carrying a
benign/malignant label --- on which the proposed model commits $16$ errors
(Section~\ref{subsec:perclass});
the $0.30$ percentage-point gap to the leader is therefore a difference of
roughly five images, well inside the noise band of a single training run. We
therefore do not claim a binary-task improvement; the proposed model is
presented as binary-task \emph{parity} with the leader, not as a binary-task
advance.

A single baseline visibly under-performs. ConvNeXt-Tiny reaches only
$95.32\,\%$ binary and $90.46\,\%$ subtype accuracy, consistent with the
well-documented data-hunger of ConvNeXt-style architectures~\cite{liu2022convnext}
when trained from scratch on small medical datasets; it consumed the full
$200$-epoch budget without satisfying the early-stopping criterion. The row is
retained rather than discarded because it documents a real and reproducible
limitation of from-scratch ConvNeXt training under this regime, and because a
benchmark that silently omits its weakest entrant misrepresents the difficulty
of the task.

\subsection{Per-Class Subtype Performance}\label{subsec:perclass}

Aggregate subtype accuracy can mask uneven behaviour across the seven classes,
so Table~\ref{tab:perclass} reports per-class precision, recall, and F1 (in
that order) for the proposed model and the two strongest baselines on the
subtype head. Five of the seven classes reach a perfect $1.00$ F1 for the
proposed model; the residual error is confined to MC and OC, two malignant
carcinomas whose visual presentations overlap.

\begin{table}[!t]
\caption{Per-class subtype performance on the test set, reported as
precision~/ recall~/ F1. Per-class support is given in parentheses. Five of
seven classes reach $1.00$ F1 for the proposed model.}
\label{tab:perclass}
\centering
\resizebox{\columnwidth}{!}{%
\begin{tabular}{l c c c}
\hline
Class & Custom V2 (Hub v2) & EfficientNet V2-B2 & Inception V3 \\
\hline
CaS (256) & \textbf{1.00 / 1.00 / 1.00} & 1.00 / 0.99 / 0.99 & 0.99 / 1.00 / 1.00 \\
CoS (239) & \textbf{1.00 / 1.00 / 1.00} & 0.99 / 1.00 / 0.99 & 1.00 / 1.00 / 1.00 \\
Gum (192) & \textbf{0.99 / 1.00 / 1.00} & 0.99 / 1.00 / 1.00 & 1.00 / 1.00 / 1.00 \\
MC (288)  & \textbf{1.00 / 0.98 / 0.99} & 1.00 / 0.99 / 0.99 & 1.00 / 0.98 / 0.99 \\
OC (173)  & \textbf{0.98 / 0.99 / 0.99} & 0.98 / 0.99 / 0.99 & 0.97 / 0.99 / 0.98 \\
OLP (288) & \textbf{1.00 / 1.00 / 1.00} & 1.00 / 0.99 / 0.99 & 1.00 / 0.99 / 0.99 \\
OT (210)  & \textbf{1.00 / 1.00 / 1.00} & 0.99 / 1.00 / 1.00 & 0.98 / 0.99 / 0.99 \\
\hline
\end{tabular}%
}
\end{table}

The proposed model's macro-averaged F1 is $0.99$, with no class falling below
$0.99$ --- a uniformity that neither of the two strongest baselines exhibits.
EfficientNet~V2-B2 and Inception~V3 each let a class slip to $0.98$ on at least
one of the three reported metrics, whereas the proposed model holds the floor
at $0.99$ across all seven classes.

To establish that this behaviour is a property of the data rather than of any
single model, Table~\ref{tab:perclass_full} reports the full per-class
precision~/ recall~/ F1 for every baseline together with the proposed model on
the subtype head. Two clinically meaningful observations follow. First, the MC
and OC columns are the only ones in which any model loses F1 below $1.00$, and
the \emph{same} two columns are where every model loses --- the residual
confusion is therefore data-driven rather than model-driven, a direct
consequence of MC and OC both being malignant carcinomas with overlapping
visual presentation. Second, the proposed model recovers a perfect
$1.00 / 1.00 / 1.00$ on five of the seven classes, a coverage no baseline
matches.

\begin{table*}[!t]
\caption{Full per-class subtype precision~/ recall~/ F1 on the test set, all
models. Per-class support is given in parentheses. The proposed model is the
only configuration whose CaS, CoS, OLP, and OT columns are all
$1.00 / 1.00 / 1.00$; residual confusion is confined to the MC and OC columns
for every model.}
\label{tab:perclass_full}
\centering
\setlength{\tabcolsep}{4pt}
\resizebox{\textwidth}{!}{%
\begin{tabular}{l c c c c c c c}
\hline
Model & CaS (256) & CoS (239) & Gum (192) & MC (288) & OC (173) & OLP (288) & OT (210) \\
\hline
ResNet50           & 1.00/0.99/1.00 & 1.00/1.00/1.00 & 0.99/0.99/0.99 & 0.98/0.98/0.98 & 0.98/0.98/0.98 & 0.99/0.99/0.99 & 1.00/1.00/1.00 \\
DenseNet121        & 1.00/0.98/0.99 & 1.00/1.00/1.00 & 1.00/1.00/1.00 & 1.00/0.99/0.99 & 0.99/0.98/0.99 & 0.98/1.00/0.99 & 1.00/1.00/1.00 \\
EfficientNet-B0    & 1.00/0.98/0.99 & 0.99/1.00/0.99 & 0.99/1.00/0.99 & 0.99/0.99/0.99 & 0.98/0.98/0.98 & 0.98/1.00/0.99 & 1.00/0.98/0.99 \\
EfficientNet V2-B2 & 1.00/0.99/0.99 & 0.99/1.00/0.99 & 0.99/1.00/1.00 & 1.00/0.99/0.99 & 0.98/0.99/0.99 & 1.00/0.99/0.99 & 0.99/1.00/1.00 \\
Inception V3       & 0.99/1.00/1.00 & 1.00/1.00/1.00 & 1.00/1.00/1.00 & 1.00/0.98/0.99 & 0.97/0.99/0.98 & 1.00/0.99/0.99 & 0.98/0.99/0.99 \\
Custom V2 (Hub v1, full) & 0.98/0.99/0.99 & 0.99/1.00/1.00 & 1.00/1.00/1.00 & 0.99/0.99/0.99 & 0.99/0.97/0.98 & 0.99/1.00/0.99 & 1.00/0.99/0.99 \\
\hline
\textbf{Custom V2 (Hub v2)} & \textbf{1.00/1.00/1.00} & \textbf{1.00/1.00/1.00} & \textbf{0.99/1.00/1.00} & \textbf{1.00/0.98/0.99} & \textbf{0.98/0.99/0.99} & \textbf{1.00/1.00/1.00} & \textbf{1.00/1.00/1.00} \\
\hline
\end{tabular}%
}
\end{table*}

The proposed model is the only configuration whose CaS, CoS, OLP, and OT
columns are all $1.00 / 1.00 / 1.00$. The residual confusion costs the model
exactly two misclassified MC images and one misclassified OC image, out of
supports of $288$ and $173$ respectively. Figure~\ref{fig:heatmap} renders the
same per-class F1 values as a model~$\times$~class heat map; what the heat map
adds over Table~\ref{tab:perclass_full} is the visual gestalt --- the MC and OC
columns form a single coherent cool band spanning every one of the ten models
at once, making the data-driven nature of the residual confusion legible at a
glance in a way a tabulated grid of numbers does not, while the proposed
model's row is the most uniformly saturated. Figure~\ref{fig:confusion}
presents the binary and subtype confusion matrices for the proposed model. The
subtype matrix is computed over the $1\,646$-image DS2 test set and the binary
matrix over the $1\,711$-image binary test set, so that both denominators are
verifiable. The binary matrix shows $11$ benign images predicted malignant and
$5$ malignant images predicted benign --- $16$ errors in $1\,711$ images,
consistent with the $99.06\,\%$ binary accuracy reported in
Table~\ref{tab:benchmark} --- and the subtype matrix confirms that every
off-diagonal entry is concentrated in the MC and OC rows, with all other
classes perfectly resolved.

\begin{figure*}[!t]
\centering
\includegraphics[width=\textwidth]{fig04_per_class_f1_heatmap}
\caption{Per-class F1 score across models, rendered as a
model~$\times$~subtype-class heat map. Darker cells indicate higher F1. The MC
and OC columns form the coolest band for every model, confirming that the
residual subtype confusion is data-driven; the proposed Custom EfficientNet~V2
(Hub~v2) row, outlined, is the most uniformly saturated.}
\label{fig:heatmap}
\end{figure*}

\begin{figure}[!t]
\centering
\includegraphics[width=\columnwidth]{fig04c_confusion_matrices_proposed}
\caption{Binary (left, $n=1\,711$) and subtype (right, $n=1\,646$) confusion
matrices for the proposed Custom EfficientNet~V2 (Hub~v2) on the held-out test
set. Off-diagonal mass in the subtype matrix is confined almost entirely to the
MC and OC rows; the remaining five classes are resolved without error.}
\label{fig:confusion}
\end{figure}

\subsection{Efficiency Comparison}\label{subsec:efficiency}

Beyond the static parameter and GFLOPs counts of Table~\ref{tab:benchmark},
deployment cost is governed by inference latency, peak GPU memory, and
training time. Table~\ref{tab:efficiency} reports these for all ten models.
Inference latency is measured at batch size~1 on the RTX~4060~Ti, after a
$200$-iteration warm-up and over $1\,000$ timed iterations; the median (P50)
and $95$th-percentile (P95) latencies are tabulated together with peak GPU
memory during inference, total wall-clock training time, and the number of
epochs completed before early stopping.

\begin{table}[!t]
\caption{Inference and training efficiency. Latency is measured at batch
size~1 on an NVIDIA RTX~4060~Ti. Bold marks the best value in each
column for the proposed model.}
\label{tab:efficiency}
\centering
\resizebox{\columnwidth}{!}{%
\begin{tabular}{l d{4} d{4} d{5} d{5} d{3}}
\hline
\multicolumn{1}{l}{Model} &
\multicolumn{1}{c}{P50 lat (ms)} &
\multicolumn{1}{c}{P95 lat (ms)} &
\multicolumn{1}{c}{GPU peak (MB)} &
\multicolumn{1}{c}{Train (min)} &
\multicolumn{1}{c}{Epochs} \\
\hline
ResNet50           &  6.88 &  9.18 & 148.37 & 131.82 & 129 \\
DenseNet121        & 17.59 & 23.77 &  80.19 & 108.84 & 100 \\
ConvNeXt-Tiny      &  6.56 &  8.80 & 171.98 & 262.55 & 200 \\
Swin-T             & 11.54 & 12.18 & 179.33 & 145.52 & 118 \\
EfficientNet-B0    &  8.88 &  9.51 &  76.91 & 104.05 & 126 \\
EfficientNet V2-B2 & 13.15 & 14.02 &  72.47 & 114.20 & 135 \\
EfficientNet V2-B3 & 15.41 & 20.38 &  94.14 &  92.47 &  99 \\
EfficientNet V2-S  & 19.12 & 25.60 & 126.94 &  98.72 &  86 \\
Inception V3       & 17.57 & 17.76 & 132.18 & 135.22 & 153 \\
\hline
\multicolumn{1}{l}{\textbf{Custom V2 (Hub v1)}} &
10.16 & 13.41 & 52.32 & 131.70 & 172 \\
\multicolumn{1}{l}{\textbf{Custom V2 (Hub v2)}} &
9.86 & 15.50 &
\multicolumn{1}{c}{\textbf{52.27}} &
\multicolumn{1}{c}{\textbf{30.80}} &
141 \\
\hline
\end{tabular}%
}
\end{table}

Two efficiency advantages stand out. The proposed model is the most
memory-efficient architecture on the GPU, with a peak footprint of
$52.27$~MB during inference --- between $30\,\%$ and $70\,\%$ lower than the
comparators, and a direct consequence of the trimmed five-stage backbone and
the lightweight AttentionHub replacing the donor's heavier Block-4. The second
advantage is convergence speed: under the identical recipe, Hub~v2 reaches its
early-stopping criterion in $141$ epochs against $172$ for Hub~v1, a $0.82\times$
epoch count consistent with the lighter, sequential topology --- a single
Triplet~\arrowto{}SE cascade imposes a smaller optimisation burden than v1's
three parallel branches followed by a fusion projection.

The wall-clock training times of Table~\ref{tab:efficiency} ($30.80$ versus
$131.70$~min) point in the same direction, but they should be read as
indicative rather than as a controlled measurement. All runs share a nominally
identical recipe and the two custom-model variants differ by only $0.4\,\%$ in
GFLOPs, yet the per-epoch wall-clock times differ by roughly $3.5\times$; this
gap far exceeds what the FLOPs difference can explain. Because training was
performed on a shared workstation without controlled process isolation,
background load can perturb per-epoch wall-clock time between runs, and we do
not attribute the full $3.5\times$ per-epoch gap to architecture. We therefore
report training time faithfully from the per-run logs but draw no efficiency
headline from it; the $141$-versus-$172$-epoch convergence difference is the
more controlled observation. Inference latency is mid-pack and competitive with
the EfficientNet family, the expected regime for a model whose GFLOPs are the
lowest in the benchmark; the proposed model is not the fastest per image, but it
is the lightest to host.
